\section{Interpret skip and changed counts in pixels}
A delta payload is a sequence of commands. Each command starts with a two-byte skip
count and a two-byte changed count, followed by three XOR bytes for each changed
pixel. The counts are big-endian; the XOR triples are red, green and blue. Start the
pixel cursor at zero for every delta record.

A skip advances the cursor without modifying those pixels. A changed count consumes
triples and XORs their numeric word values into the following pixels. Both counts
may be zero individually, but their sum must be positive. Require each count to fit
the pixels still available, using the remaining count after the skip when admitting
the changed span. Require every changed triple to fit the remaining payload bytes.
Stop only when all 4096 pixels are covered, then require exact payload completion.

A final command may consist solely of a nonzero skip. A final changed span needs no
additional terminator. Do not append a zero/zero command: it makes no progress and
would leave the cursor unchanged. A changed span is allowed to include a zero triple.
That is a valid, potentially inefficient way to encode an unchanged pixel. Accept
such equivalent encodings when grading decoded results.

There can be at most 4096 commands because every command covers at least one pixel.
Their four-byte headers and at most 4096 triples give the conservative payload bound
$4\cdot4096+3\cdot4096=28\,672$. This bounds storage; it is not a required encoded size.
